% Project: `` Anchor institutions as macroeconomic stabilizers"
% Section: Conclusion

%\newpage
\section{ Conclusion}

\noindent This project examines how university-dependent local economies respond to national recessions and whether the stabilizing influence of universities persists when campus activity is disrupted. The preliminary evidence, based on an external definition of college towns and unemployment data from counties and metropolitan areas, suggests two broad patterns. First, college-town areas tend to exhibit lower unemployment rates in normal periods and experience similar or slightly smaller unemployment increases in the 2001 and 2008 recessions. Second, during the COVID recession, when in-person university activity was severely curtailed, these differences largely disappear, indicating that the stabilizing channel may weaken when the campus itself becomes constrained. Although these findings are exploratory, they are consistent with the broader stability and fragility mechanisms highlighted in the literature.

The next stage of the project will replace the external classification with an endogenous measure of university exposure constructed from QCEW and IPEDS data. This will allow for more precise identification of college towns, continuous exposure measures, and richer heterogeneity analyses. With the full dataset in place, the empirical design will combine synthetic difference in differences, interaction difference in differences, and local projections to estimate both average and dynamic effects of university dependence across recessions. A separate event-study will focus on the 2020 campus shutdown, and further extensions will incorporate sectoral employment, payroll, and housing market outcomes, as well as spillover effects on neighboring areas.

In terms of feasibility, the remaining data work is substantial but attainable: QCEW extraction is labor intensive, yet all needed sources are available, and the merging process is straightforward once compiled. The empirical strategies are well established and suitable for the research question. By placing universities within the broader question of how anchor institutions shape regional resilience, this project provides new evidence on when stable local anchors act as buffers and when they become points of vulnerability. It also expands the focus beyond flagship institutions to a more comprehensive set of college towns, offers a dynamic view of adjustment across multiple recessions, and integrates labor market and housing outcomes.






